% !TEX TS-program = pdflatex
% !TeX program = pdflatex
% !TEX encoding = UTF-8
% !TEX spellcheck = fr

\documentclass[xcolor=table]{beamer}


%\usepackage{fullpage}
%\usepackage[left=2.8cm,right=2.2cm,top=2 cm,bottom=2 cm]{geometry}
\setbeamersize{text margin left=10pt,text margin right=10pt}


\usepackage{amsmath,amssymb} 
\usepackage{fontenc}
\usepackage{textcomp}
\usepackage[utf8]{inputenc}
\usepackage{CJKutf8}
\usepackage[french]{babel}
\usepackage{arabtex,acolor}
\usepackage{txfonts}
\usepackage{tipa}
\usepackage[]{graphicx}
\usepackage{multirow}
\usepackage{hyperref}
\usepackage{colortbl}
\usepackage{tabularray}
\usepackage{listingsutf8}
\usepackage{wrapfig}
\usepackage{multicol}
\usepackage[export]{adjustbox} %for images in table, also for frame
\usepackage[many]{tcolorbox}
\usepackage{wasysym}
\usepackage[french,lined]{algorithm2e}
\usepackage{alltt} %verbatim with commands
\usepackage{longtable}
\usepackage{tabu}


\definecolor{lightblue}{HTML}{D0D2FF}
\definecolor{lightyellow}{HTML}{FFFFAA}
\definecolor{darkblue}{HTML}{0000BB}
\definecolor{olivegreen}{HTML}{006600}
\definecolor{darkgreen}{HTML}{008B45} %009B55
\definecolor{violet}{HTML}{6600CC}
%\definecolor{deeppink}{HTML}{FF1493}
\definecolor{orangey}{HTML}{FFBB00}

\usetheme{Warsaw} % Antibes Boadilla Warsaw

\beamertemplatenavigationsymbolsempty


\let\oldcite\cite
\renewcommand{\cite}[1]{{\bfseries\color{orangey}\oldcite{#1}}}

\tcbuselibrary{listings}

\newcommand{\kurl}[1]{{\scriptsize\bfseries\color{orangey}\url{#1}}}


%\renewcommand{\baselinestretch}{1.5}

\def\supit#1{\raisebox{0.8ex}{\small\it #1}\hspace{0.05em}}

\AtBeginSection{%
	\begin{frame}
		\sectionpage
	\end{frame}
}

\newcommand{\rottext}[2]{%
	\rotatebox{90}{%
	\begin{minipage}{#1}%
		\raggedleft#2%
	\end{minipage}%
	}%
}


\institute{ %
Laboratoire de la Communication dans les Systèmes Informatiques (LCSI)

\vspace{6pt}
École  nationale Supérieure d'Informatique (ESI, ex. INI), Alger, Algérie
}
\author[ \textbf{\footnotesize\insertframenumber/\inserttotalframenumber} \hspace*{1cm} ESI - ARIES Abdelkrime (Master 2022/2023)] %
{ARIES Abdelkrime}
%\titlegraphic{\includegraphics[height=1cm]{../img/esi-logo.png}%\hspace*{4.75cm}~
	
	
	\date{Master 2022/2023} %\today

\titlegraphic{%
	\includegraphics[height=1cm]{../img/esi-logo.png}%
	\hspace{2cm}
	\includegraphics[height=1cm]{../img/lcsi-logo.png}%
}


%\setbeamertemplate{headline}{}



\newcommand{\keyword}[1]{\textcolor{red}{\bfseries\itshape #1}}
\newcommand{\expword}[1]{\textcolor{olivegreen}{#1}}
\newcommand{\optword}[1]{\textcolor{violet}{\bfseries #1}}

\makeatletter
\newcommand\mysphere{%
	\parbox[t]{10pt}{\raisebox{0.2pt}{\beamer@usesphere{item projected}{bigsphere}}}}
\makeatother

%\let\oldtabular\tabular
%\let\endoldtabular\endtabular
%\renewenvironment{tabular}{\rowcolors{2}{white}{lightblue}\oldtabular\rowcolor{blue}}{\endoldtabular}


\NoAutoSpacing %french autospacing after ":"

\def\graphpath{}

\newcommand{\changegraphpath}[1]{\def\graphpath{#1}}


\newcommand{\vgraphpage}[2][.84\textheight]{%
%	\begin{center}%
		\includegraphics[height=#1]{\graphpath #2}%
%	\end{center}%
}

\newcommand{\hgraphpage}[2][\textwidth]{%
%	\begin{center}%
		\includegraphics[width=#1]{\graphpath #2}%
%	\end{center}%
}

\newcommand{\graphpage}[2][]{%
	\includegraphics[#1]{\graphpath #2}%
}

\bibliographystyle{apalike}

\newcommand{\insertbibliography}[2]{
	\appendix
	\section*{Bibliographie}
	\nocite{#2}
%	\makeatletter % to change template
%	\setbeamertemplate{headline}[default] % not mandatory, but I though it was better to set it blank
%	\def\beamer@entrycode{\vspace*{-\headheight}} % here is the part we are interested in :)
%	\makeatother
	\begin{multicols*}{2}[\frametitle{\insertsection} \usebeamertemplate*{frametitle}]%\usebeamertemplate*{frametitle}\frametitle{Références}
		\tiny
		\bibliography{#1}
	\end{multicols*}
}

\definecolor{my-grey}{RGB}{233, 233, 233}

\newcommand{\insertlicence}{
	\begin{frame}[plain]
	\frametitle{Licence : CC-BY 4.0}
%	\framesubtitle{Licence: CC-BY-NC 4.0}

	\begin{tcolorbox}[colback=cyan,
		colframe=cyan,  
		arc=0pt,outer arc=0pt,
		valign=top, 
		halign=center,
		width=\textwidth]
		
		\includegraphics[width=.5cm]{../img/licence/cc_icon_white_x2.png}
		\includegraphics[width=.5cm]{../img/licence/attribution_icon_white_x2.png}
		
		\color{white}
		\bfseries Attribution 4.0 International (CC BY 4.0) \\
		\tiny \url{https://creativecommons.org/licenses/by/4.0/deed.fr}
		
	\end{tcolorbox}\vspace{-.5cm}
	\begin{tcolorbox}[colback=my-grey,
		colframe=my-grey,  
		center, arc=0pt,outer arc=0pt,
		valign=top, 
		halign=left,
		width=\textwidth]
		
		\tiny
		
		\begin{center}
			\bfseries\Large
			Vous êtes autorisé à :
		\end{center}
		
		\begin{minipage}{0.83\textwidth}
			\begin{itemize}
				\item[] \textbf{Partager} — copier, distribuer et communiquer le matériel par tous moyens et sous tous formats
				\item[] \textbf{Adapter} — remixer, transformer et créer à partir du matériel
				pour toute utilisation, y compris commerciale.
			\end{itemize}
		\end{minipage}
		\begin{minipage}{0.15\textwidth}
			\includegraphics[width=\textwidth]{../img/licence/FreeCulturalWorks_seal_x2.jpg}
		\end{minipage}
	
		
		\begin{center}
			\bfseries\Large
			Selon les conditions suivantes :
		\end{center}
		
		\begin{itemize}
			\item[] \textbf{Attribution} — Vous devez créditer l'Œuvre, intégrer un lien vers la licence et indiquer si des modifications ont été effectuées à l'Oeuvre. Vous devez indiquer ces informations par tous les moyens raisonnables, sans toutefois suggérer que l'Offrant vous soutient ou soutient la façon dont vous avez utilisé son Oeuvre. 
			\item[] \textbf{Pas de restrictions complémentaires} — Vous n'êtes pas autorisé à appliquer des conditions légales ou des mesures techniques qui restreindraient légalement autrui à utiliser l'Oeuvre dans les conditions décrites par la licence.
		\end{itemize}
		
	\end{tcolorbox}
	
%	\begin{center}
%		\bfseries Attribution 4.0 International (CC BY 4.0)
%		\url{https://creativecommons.org/licenses/by/4.0/deed.fr}
%	\end{center}

%	\tiny
%
%	Vous êtes autorisé à : 
%	\begin{itemize}
%		\item \textbf{Partager} — copier, distribuer et communiquer le matériel par tous moyens et sous tous formats
%		\item \textbf{Adapter} — remixer, transformer et créer à partir du matériel
%	\end{itemize}
%	
%	Selon les conditions suivantes : 
%	\begin{itemize}
%		\item \textbf{Attribution} — Vous devez créditer l'Œuvre, intégrer un lien vers la licence et indiquer si des modifications ont été effectuées à l'Oeuvre. Vous devez indiquer ces informations par tous les moyens raisonnables, sans toutefois suggérer que l'Offrant vous soutient ou soutient la façon dont vous avez utilisé son Oeuvre.
%		\item \textbf{Pas d'Utilisation Commerciale} — Vous n'êtes pas autorisé à faire un usage commercial de cette Oeuvre, tout ou partie du matériel la composant. 
%		\item \textbf{Pas de restrictions complémentaires} — Vous n'êtes pas autorisé à appliquer des conditions légales ou des mesures techniques qui restreindraient légalement autrui à utiliser l'Oeuvre dans les conditions décrites par la licence.
%	\end{itemize}

	\end{frame}
}

\settowidth{\leftmargini}{\usebeamertemplate{itemize item}}
\addtolength{\leftmargini}{\labelsep}

\AtBeginDocument{
	\newcolumntype{L}[2]{>{\vbox to #2\bgroup\vfill\flushleft}p{#1}<{\egroup}} 
	
	\begin{frame}[plain]
		\maketitle
	\end{frame}

	\insertlicence
}


% needs etoolbox; to break links after -
\appto\UrlBreaks{\do\-}


\makeatletter
\newcommand{\xRightarrow}[2][]{\ext@arrow 0359\Rightarrowfill@{#1}{#2}}
\makeatother


\usefonttheme{structurebold}
%\usefonttheme{professionalfonts}

\title[TALN : 05- Analyse syntaxique]%
{Traitement automatique du langage naturel\\Chapitre 05 : Analyse syntaxique} 

\changegraphpath{../img/syntaxe/}

\begin{document}
	
\begin{frame}
\frametitle{Traitement automatique du langage naturel}
\framesubtitle{Analyse syntaxique : Introduction}

\begin{exampleblock}{Exemple d'une phrase en français}
	\begin{center}
		\Large\bfseries
		L'élève a écrit une solution avec un stylo. 
		
		L'élève a écrit une solution avec une explication.
	\end{center}
\end{exampleblock}

\begin{itemize}
	\item Qui a écrit la solution ?
	\item Quel instrument a-t-il utilisé pour l'écriture ?
	\begin{itemize}
		\item Un stylo ? Une explication ?
	\end{itemize}
	\item Est-ce qu'il y avait d'autre chose en plus de la solution ?
	\begin{itemize}
		\item Un stylo ? Une explication ?
	\end{itemize}
	\item Comment avez-vous déduit ça ?
\end{itemize}

\end{frame}

%\begin{frame}
%\frametitle{Traitement automatique du langage naturel}
%\framesubtitle{Analyse syntaxique : Un peu d'humour}
%
%\begin{center}
%	\vgraphpage{humour-pos.jpg}
%\end{center}
%
%\end{frame}

\begin{frame}
\frametitle{Traitement automatique du langage naturel}
\framesubtitle{Analyse syntaxique : Plan}

\begin{multicols}{2}
%	\small
\tableofcontents
\end{multicols}
\end{frame}

%===================================================================================
\section{Structures syntaxiques}
%===================================================================================

\begin{frame}
\frametitle{Analyse syntaxique}
\framesubtitle{Structures syntaxiques}

\begin{itemize}
	\item 
\end{itemize}

\end{frame}

\subsection{Grammaires à contexte libre}

\begin{frame}
\frametitle{Analyse syntaxique : Structures syntaxiques}
\framesubtitle{Grammaires à contexte libre}

\begin{itemize}
	\item $G <X,V, P, S>$ est une grammaire.
	\item $X$ est l'alphabet. Dans notre cas, c'est l'ensemble des mots.
	\begin{itemize}
		\item \expword{X = \{le, petit, chat, mange, un, poisson, ...\}}
	\end{itemize}
	\item $V$ est l'ensemble des  variables.
	\begin{itemize}
		\item \expword{V = \{S, NP, VP, DET, N, ADJ, ...\}}
	\end{itemize}
	\item $S \in V$ est l'axiome.
	\item $P$ est l'ensemble des règles de production.
	\item Les règles sont de la forme $A \rightarrow \beta \text{ avec } A \in V,\, \beta \in (X \cup V)^*$
	\begin{itemize}
		\item \expword{S \textrightarrow NP VP}
		\item \expword{NP \textrightarrow DET ADJ N \textbar\ DET N}
		\item \expword{VP \textrightarrow V NP}
	\end{itemize}
\end{itemize}

\end{frame}

\subsection{Annotation constituante}

\begin{frame}
\frametitle{Analyse syntaxique : Structures syntaxiques}
\framesubtitle{Annotation constituante}

\begin{itemize}
	\item 
\end{itemize}

\end{frame}

\subsection{Annotation fonctionnelle}

\begin{frame}
\frametitle{Analyse syntaxique : Structures syntaxiques}
\framesubtitle{Annotation fonctionnelle}

\begin{itemize}
	\item 
\end{itemize}

\end{frame}

%===================================================================================
\section{Analyse des constituants}
%===================================================================================

\begin{frame}
\frametitle{Analyse syntaxique}
\framesubtitle{Analyse des constituants}

%\begin{itemize}
%	\item Méthodes classiques 
%	\begin{itemize}
%		\item CKY
%		\item LR 
%	\end{itemize}
%	\item Apprentissage automatique 
%	\begin{itemize}
%		\item CKY probabiliste
%	\end{itemize}
%\end{itemize}

\begin{tabular}{|p{.3\textwidth}|p{.3\textwidth}|p{.3\textwidth}|}
	\hline
	& Méthodes classiques & Apprentissage automatique \\
	\hline
	Analyse ascendante & CKY & CKY probabiliste\\
	\hline
	Analyse descendante &  & \\
	\hline
\end{tabular}

\end{frame}

\subsection{Algorithme CKY}

\begin{frame}
\frametitle{Analyse syntaxique : Analyse des constituants}
\framesubtitle{Algorithme CKY (1)}

\begin{itemize}
	\item Algorithme de Cocke-Younger-Kasami
	\item Analyse ascendante
	\item Prérequis : Forme normale de Chomsky
\end{itemize}

\begin{definition}[FNC : Forme normale de Chomsky]
	\[
	A \rightarrow  B C \text{ où } A, B, C \in V
	\]
	
	\[
	A \rightarrow a \text{ où } a \in X
	\]
\end{definition}

\end{frame}

\begin{frame}
\frametitle{Analyse syntaxique : Analyse des constituants}
\framesubtitle{Algorithme CKY (2)}

\begin{block}{CKY}
	\scriptsize
	\begin{algorithm}[H]
		\KwData{une grammaire $G <X,V, P, S>$ en FNC, une phrase $w = w_1 \ldots w_n$}
		\KwResult{$meilleur\_chemin$, $prob\_chemin$}
		
%		Créer une matrice $viterbi[N, T]$\;
%		
%		\Pour{état $ s = 1 \ldots N$}{
%			$viterbi[s, 1] = \pi_s * b_s(w_1);\, backpointer[s, 1] = 0$ \;
%		}
%		
%		\Pour{$ t = 1 \ldots T$}{
%			\Pour{état $ s = 1 \ldots N$}{
%				$viterbi[s, t] = \max\limits_{s'=1}^N viterbi[s', t-1] * a_{s',s} * b_s(w_t)$\;
%				$backpointer[s, t] = \arg\max\limits_{s'=1}^N viterbi[s', t-1] * a_{s',s} * b_s(w_t)$\;
%			}
%		}
%		
%		$prob\_chemin = \max\limits_{s=1}^N viterbi[s, T];\, pointeur\_chemin = \arg\max\limits_{s=1}^N viterbi[s, T]$\;
%		
%		$meilleur\_chemin$ est le chemin qui commence par $pointeur\_chemin$ et qui suit $backpointer$
%		
%		\Return $meilleur\_chemin$, $prob\_chemin$\;
		
	\end{algorithm}
\end{block}

\end{frame}

\subsection{Algorithme CKY probabiliste}

\begin{frame}
\frametitle{Analyse syntaxique : Analyse des constituants}
\framesubtitle{Algorithme CKY probabiliste}

\begin{itemize}
	\item 
\end{itemize}

\end{frame}

\subsection{Algorithme LR}

\begin{frame}
\frametitle{Analyse syntaxique : Analyse des constituants}
\framesubtitle{Algorithme LR}

\begin{itemize}
	\item 
\end{itemize}

\end{frame}

%===================================================================================
\section{Analyse des dépendances}
%===================================================================================

\begin{frame}
\frametitle{Analyse syntaxique}
\framesubtitle{Analyse des dépendances}

\begin{itemize}
	\item 
\end{itemize}

\end{frame}

\subsection{shift-reduce}

\begin{frame}
\frametitle{Analyse syntaxique : Analyse des dépendances}
\framesubtitle{shift-reduce}

\begin{itemize}
	\item 
\end{itemize}

\end{frame}

\subsection{Arc-Eager}

\begin{frame}
\frametitle{Analyse syntaxique : Analyse des dépendances}
\framesubtitle{Arc-Eager}

\begin{itemize}
	\item 
\end{itemize}

\end{frame}

\subsection{Arbre couvrant de poids minimal}

\begin{frame}
\frametitle{Analyse syntaxique : Analyse des dépendances}
\framesubtitle{Arbre couvrant de poids minimal}

\begin{itemize}
	\item 
\end{itemize}

\end{frame}

\insertbibliography{TALN05}{*}

\end{document}

